\chapter{Source concordance, review coverage, and mathematical provenance}
\label{ch:source-concordance}

\section{What the review establishes}

The mathematical argument in this monograph rests on several different kinds
of evidence. A bibliographic inventory identifies sources. Text extraction
makes a document searchable. Reading a passage identifies its claims and
hypotheses. Reconstructing a proof checks its inference. A successful Lean
compilation checks a precisely stated formal declaration relative to its
dependencies. These activities support different conclusions, and this
chapter records their actual extent. The resulting account is substantial
but is not a claim to have proved, formalized, or independently verified every
statement in every linked document.

There are two relevant authorship layers. The supplied classical books on
lattices over orders are by Klaus W. Roggenkamp, with Verena Huber-Dyson as
coauthor of Volume~I; Volume~II is authored by Roggenkamp
\cite{RoggenkampHuberDyson1970,v2Roggenkamp1970II}. The modern research
repositories and linked surface-order manuscripts are associated with
Amelie Schreiber \cite{GRHRepo,SchreiberSurfaceI,SchreiberSurfaceII}.
Neither author's complete research oeuvre has been reviewed. Our use of
selected additional papers by Roggenkamp concerns concrete questions about
orders, irreducible lattices, and Bass-type phenomena
\cite{Roggenkamp1969Irreducible,Roggenkamp1969Separable,v2RoggenkampBass1971}.
It does not license a statement about everything either author has written.

\section{The fixed source snapshot}

The primary repository was examined at commit
\texttt{dbd954eede\allowbreak{}479a512983\allowbreak{}1e4e7c3053\allowbreak{}7389ae973d}
\cite{GRHRepo}. Four important source files received complete text review:
\path{book/main.tex}, the generated chapter
\path{book/chapters/reviewed-mathematics.tex}, and
\path{papers/core.tex} and \path{papers/details.tex}. Their measured
coverage is as follows. Words are whitespace-delimited source tokens, so the
counts include mathematical notation and LaTeX commands; they are coverage
measurements rather than estimates of polished prose length.

\begin{center}
\begin{tabular}{lrr}
\hline
Source file & Lines & Source words\\
\hline
Book entry point & 67 & 466\\
Reviewed mathematics chapter & 1453 & 23652\\
Core paper & 1072 & 6123\\
Details paper & 1621 & 9893\\
\hline
Total & 4213 & 40134\\
\hline
\end{tabular}
\end{center}

Complete reading of these files does not mean that every transitive
\texttt{input} was reviewed in the same audit. The book entry point includes
many other chapters. The two papers also include separate short-row-pair and
M\"obius-splitting files. Each source retains its own coverage status in the
companion review. Likewise, reading a discussion of an external theorem
does not constitute an independent proof audit of that theorem's original
paper. The distinction matters especially for analytic number theory,
operator algebras, derived equivalences, and arithmetic descent.

The linked mathematics submodule was inventoried at commit
{\small\texttt{4eaee87548\allowbreak{}250b2682af\allowbreak{}fda7d05d94\allowbreak{}7c05554bae}}. Its tree contains
174,524 files. This is a file-tree inventory, not a claim that those files
were all read. The repository link inventory contains 409 distinct external
URLs. A URL may identify a paper, a documentation page, a repository, or a
repeated discussion of a result; it is not automatically one independent
theorem or one reviewed document. These 409 destinations have not all been
reviewed. The two linked author manuscripts, with 223 and 265 PDF pages at
their pinned snapshots, were extracted and selectively reviewed
\cite{SchreiberSurfaceI,SchreiberSurfaceII}. Their complete proofs were not
audited page by page. These explicit limits prevent a large source manifest
from being mistaken for a completed universal verification.

\section{The classical books: extraction and passage coverage}

All 309 PDF pages of Volume~I were extracted by optical character recognition,
and all resulting page text was read. Selected mathematical statements were
also checked against page images. Formula reconstruction remains selective:
OCR can lose subscripts, reverse inclusion signs, merge displayed matrices,
or turn a quantified condition into an apparently unconditional statement.
Consequently ``all page text read'' is stronger than a table-of-contents
survey but weaker than a verified transcription or a complete proof audit
\cite{RoggenkampHuberDyson1970}.

All 392 PDF pages of Volume~II were likewise processed by OCR, and all
resulting page text has now been read across the review team. The initial
review covered 158 pages, including the complete extracted prose of
Chapters~VII and VIII, selected passages in the remaining chapters, and the
final errata. Two complementary supplemental reads covered the remaining
101 and 133 pages without gaps. The extraction contains approximately
99,149 whitespace-delimited words. It is a working aid, not an authoritative
mathematical edition. The page coverage manifests distinguish extraction,
complete OCR-content reading, and selected visual verification. Some densely
typeset classifications and matrix displays remain badly garbled in OCR;
their presence and role were reviewed, but their formulas and proofs were
not thereby certified \cite{v2Roggenkamp1970II}.

Volume~II also uses two page coordinate systems. Its printed page~114 is
PDF page~119; printed page~139 is PDF page~144; and the printed errata
page~387 is PDF page~392. Keeping these coordinates explicit makes the most
important comparisons reproducible. Full copyrighted scans and bulk OCR are
not reproduced as new research content. The public companion instead
contains original notes, coverage metadata, bibliographic identifiers, and
the new derivations needed by the present argument.

\section{Three decisive anchors in Volume~II}

The first anchor is the discussion of topological connections on printed
page~114. It motivates a bridge from modules over rings of continuous
functions to vector bundles, but a modern statement must specify which
K-theory is meant. For a compact Hausdorff space, the finite projective
module--vector bundle correspondence gives the familiar interpretation of
\(K_0(C(X,\mathbb C))\) as complex topological \(K^0(X)\)
\cite{v2Roggenkamp1970II,v2Swan1962}. This does not identify algebraic
and topological K-theory in every degree. Already at a point,
\(K_1^{\mathrm{alg}}(\mathbb C)=\mathbb C^\times\), whereas
\(K_1^{\mathrm{top}}(\mathbb C)=0\). The present manuscript therefore
uses the historical connection with its category, topology, and degree
specified, rather than allowing the notation \(K_1\) alone to carry an
unjustified comparison.

The second anchor is the weighted Gram-rank argument on printed page~139.
For a real matrix \(X\) and a positive diagonal matrix \(W\),
\[
 \ker(X^{\mathsf T}WX)=\ker X,
 \qquad
 \operatorname{rank}(X^{\mathsf T}WX)=\operatorname{rank}X.
\]
Indeed, \(v^{\mathsf T}X^{\mathsf T}WXv\) is a sum of positive weights
times squares. Its vanishing forces every component of \(Xv\) to vanish.
This elementary argument, checked against the scan, supplies a genuine
connection to the Hodge quadratic forms developed in the complexes chapter
\cite{v2Roggenkamp1970II}. It also identifies the necessary hypothesis:
positivity of the coefficient form. With an indefinite weight, cancellation
can invalidate the kernel identity. Over a finite field, a nonzero vector
can be isotropic, so a real positivity argument cannot simply be transferred.
The algebraic identity and the ordered-field inference must remain separate.

The third anchor is the 1970 errata on printed page~387. It includes a
restriction involving totally definite quaternion components for a statement
in Volume~I and a correction to the argument on Volume~I, chapter-page~IV.13.
The latter concerns an unjustified order-linear splitting of torsion and
torsion-free parts. The replacement argument uses the appropriate local
and faithfully flat module-theoretic methods
\cite{v2Roggenkamp1970II,RoggenkampHuberDyson1970}. This correction is
material to any attempt to reuse the books for equivalences of module
categories. Our matrix-ring Morita calculation has an independent direct
proof and does not invoke the faulty splitting. A historical source should
be read together with its published corrections, especially when the
correction changes the mechanism of a proof.

\section{What survives scrutiny in the upstream research program}

The fully read source files contain substantial mathematical repairs and
careful negative results. They distinguish infinite orders from finite
quotients, path degree from the formal order parameter, and algebraic,
formal, Laurent, and operator-norm completions. Their nodal examples have
explicit alternating resolutions. Their semi-invariant constructions
distinguish determinant weights from ordinary invariants, and their
representation-theoretic examples state where Schur constructions cease to
preserve additivity or a chosen scalar action \cite{GRHRepo}. These are
useful structural results even when they do not settle the eventual analytic
problem.

The same source explicitly identifies limitations of several proposed
arithmetic transfers. A complex dessin does not by itself recover rational
descent data. Inserting a known two-dimensional arithmetic matrix into a
larger projective module changes an ordinary determinant by multiplicity.
A later branch coefficient functor can remove a particular multiplicity,
but it changes the determinant target and gives up specified symmetries.
Tame descent can preserve an endomorphism order while failing to preserve
the desired integral column module. None of these distinctions can be
replaced by a generic assertion that the categories are ``equivalent''
\cite{GRHRepo}.

The analytic papers are equally explicit about the missing global step.
Self-adjoint finite spectral data yield a critical-line conclusion only
after an appropriately normalized determinant converges locally uniformly
to the intended arithmetic function. Strong operator convergence alone
does not provide trace-norm determinant convergence. Finite Euler products
cannot supply a locally uniform zero-bearing limit on a region where the
approximants remain zero-free and the limit is nonzero. Positivity of an
unrelated matrix kernel does not prove positivity of an arithmetic moment
form. The papers' conditional closure statements are useful because they
locate these obligations precisely; they do not discharge them
\cite{GRHRepo}.

There is one editorial qualification worth retaining in future revisions.
An earlier discussion says that no multiplicity-removing functor has been
supplied, whereas a later chapter constructs a branch functor under changed
requirements. The earlier sentence is accurate when restricted to the naive
tensor insertion being discussed. Read globally, it is too broad. The
mathematical distinction is between failure of that insertion and success
of a narrower construction that relinquishes symmetry. The review records
this qualification rather than silently replacing either claim.

\section{From sources to the new theorems in this monograph}

The complexes chapters distinguish three kinds of dependence. Elementary
results such as the dimension of a two-step rank stratum, its tangent
equation, finite Hodge decomposition, and the orthogonal quadric calculation
receive proofs in the text. The Buchsbaum--Eisenbud exactness criterion and
the geometry of varieties of complexes are attributed to their separate
primary sources; the latter attribution includes De Concini and Strickland
\cite{v2BE1973,v2DS1981,v2DEP1982}. The Atiyah--Singer index theorem and
the analytical existence theory behind elliptic Hodge decomposition remain
external advanced results, stated with their hypotheses rather than
presented as newly proved here \cite{v2AS1968I,v2AS1968III}.

This discipline matters at the gauge-theoretic interface. The covariant
derivative squares to a curvature action. A deformation sequence is a
complex under specified background conditions, such as flatness or the
anti-self-dual equation in the appropriate projected sequence. Its index
concerns a classical elliptic problem and cannot, without another theorem,
be read as a positive spectral gap for an interacting quantum Hamiltonian
\cite{v2AHS1978,v2AB1983}. The Yang--Mills chapter separately assesses the
recent construction and its stated acceptance criteria. The source
concordance does not enlarge those claims.

\section{A finite-field correction in the source matrix reductions}

A further qualification arises in Volume~II, printed page~351
(chapter-page X~65; PDF page~356), in the reduction of extension matrices
over a cubic finite-field extension \cite{v2Roggenkamp1970II}. The scan
asserts that, outside characteristic~3, one may choose a primitive element
whose cube belongs to the base field. Visual inspection confirms that this
is an overstatement in the source, rather than an OCR error. The weaker
relation used subsequently has the following direct justification.

\begin{proposition}[Trace-zero cubic presentation]
Let \(S/k\) be an extension of finite fields of degree~3, with
\(\operatorname{char}k\ne3\). There is an element \(\alpha\) generating
\(S\) over \(k\), and elements \(a,b\in k\) with \(b\ne0\), such that
\[
 \alpha^3=a\alpha+b,
 \qquad
 \alpha^{-1}=b^{-1}\alpha^2-ab^{-1}.
\]
There need not be a generator whose cube lies in \(k\).
\end{proposition}

\begin{proof}
Choose \(\beta\in S\setminus k\). The tower formula shows that
\([k(\beta):k]\) divides~3 and is greater than~1, hence equals~3.
Thus \(\beta\) generates \(S\). The field trace is \(k\)-linear,
and the trace of \(c\in k\) is \(3c\): multiplication by \(c\) on
the three-dimensional vector space \(S\) has matrix \(cI_3\).
Division by~3 is permitted, so put
\[
 \alpha=\beta-\frac{\operatorname{Tr}_{S/k}(\beta)}3.
\]
Then \(\operatorname{Tr}_{S/k}(\alpha)=0\). Since the subtracted
quantity lies in \(k\), the element \(\alpha\) remains outside
\(k\) and generates \(S\).

Its monic minimal polynomial has degree~3. In the basis
\(1,\alpha,\alpha^2\), multiplication by \(\alpha\) has the
corresponding companion matrix. Its trace is the negative of the
quadratic coefficient, so that coefficient vanishes. We may therefore
write the polynomial as \(T^3-aT-b\). Its constant term cannot vanish:
otherwise \(T\) would divide this irreducible cubic. Consequently
\(b\ne0\), and the polynomial equation gives
\(\alpha(\alpha^2-a)=b\). Multiplication by \(b^{-1}\) proves the
inverse formula.

For the last assertion, take \(S=\mathbb F_8\) and \(k=\mathbb F_2\).
Every polynomial \(T^3-c\), with \(c\in\mathbb F_2\), has a root in
\(\mathbb F_2\): choose~0 when \(c=0\) and~1 when \(c=1\).
It is therefore reducible. It cannot be the degree-three minimal
polynomial of a generator of \(\mathbb F_8/\mathbb F_2\).
\end{proof}

This repair establishes exactly the depressed-cubic relation and quadratic
inverse expression needed at that point. It leaves the source's separate
characteristic-3 argument separate. It also illustrates the scope of the
review: correcting this local assertion does not certify every subsequent
matrix transformation or formalize the classification theorem in Lean.

\section{Formal proofs, computations, and remaining coverage}

The original Lean companion file contains 25 selected declarations compiled against
Lean~4.24.0 and the pinned mathlib revision recorded in its manifest. Their
types concern actual finite-dimensional maps, concrete matrix identities,
real positive semidefiniteness, a determinant factorization, modular kernel
relations, and noise cancellation. Each declaration has a printed dependency
report; the successful build contains no project-specific axiom or
\texttt{sorry}. Standard foundational dependencies are recorded rather than
described as an absence of all axioms. This is a precise verification result
for the declarations present in the file.

An additional seven-declaration file verifies finite Pauli identities,
binary phase and shift commutation, the chain-boundary specialization, and
a conditional integer Euler-rank substitution. Its independent successful
build and axiom audit supplied the first 32 selected declarations. The
second audit adds sixteen declarations in DiagramGrammar, bringing the
three-module inventory to 48; its scope is recorded in the formalization appendix.

It is not a formalization of every theorem in either classical book, the
entire source repository, Shor's quantum circuit, the Atiyah--Singer index
theorem, or a Yang--Mills continuum construction. Upstream reports of earlier
formal components likewise retain their literal theorem scopes. A numerical
example, a source-code formula, and a proved quantified theorem are different
deliverables even when they concern the same matrix. The companion manifests
make those distinctions visible.

The synthesis chapter's periodic order example was also checked directly:
the annihilators give the stated exact alternating resolution; tensoring
with the quotient gives the claimed odd-degree torsion groups; and the
real finite matrix specialization has middle Hodge operator
\((1+t^2)I\). Scaling the differentials then preserves middle exactness
while sending the least Hodge eigenvalue to zero. These calculations show
why the review's main limitation is mathematical as well as bibliographic:
an algebraic identification does not supply the uniform estimates required
by a spectral limit or the cost bounds required by an algorithmic claim.
Completing the unreviewed corpus would improve coverage; proving those
additional estimates would require genuinely additional mathematics.
